%% ============================================================
%% Introduction to Cognitive Psychology — Week 1 Study Notes
%% ============================================================
\documentclass[11pt,a4paper]{article}

\usepackage[a4paper, left=1.8cm, right=1.8cm, top=1.3cm, bottom=1.8cm]{geometry}
\usepackage{booktabs}
\usepackage{tabularx}
\usepackage{array}
\usepackage{enumitem}
\usepackage{titlesec}
\usepackage{fancyhdr}
\usepackage{fontenc}
\usepackage{inputenc}
\usepackage{microtype}
\usepackage{parskip}
\usepackage{hyperref}
\usepackage{amsmath}
\usepackage{amssymb}

\hypersetup{hidelinks, colorlinks=false}

%% ── Table helpers ─────────────────────────────────────────
\newcolumntype{L}[1]{>{\raggedright\arraybackslash}p{#1}}

%% ── Section headings ──────────────────────────────────────
\titleformat{\section}{\large\bfseries}{}{0em}{}[\titlerule]
\titleformat{\subsection}{\normalsize\bfseries}{}{0em}{}
\titleformat{\subsubsection}{\normalsize\bfseries}{}{0em}{}

\titlespacing{\section}{0pt}{12pt}{4pt}
\titlespacing{\subsection}{0pt}{8pt}{3pt}
\titlespacing{\subsubsection}{0pt}{6pt}{2pt}

%% ── Page style ────────────────────────────────────────────
\pagestyle{fancy}
\fancyhf{}
\renewcommand{\headrulewidth}{0pt}
\fancyfoot[C]{\small Cognitive Psychology --- Week 1 \textbar\ Page \thepage}

%% ── Misc helpers ──────────────────────────────────────────
\newcommand{\bb}[1]{\textbf{#1}}
\setlength{\parskip}{4pt}

%% ═══════════════════════════════════════════════════════════
\begin{document}
\setlength{\arrayrulewidth}{0.4pt}

\begin{center}
\bfseries\LARGE Introduction to Cognitive Psychology\par
\vspace{0.2cm}
\large Assignment 1 Detailed Solutions
\end{center}
\vspace{0.5cm}

%% ════════════════════════════════════════════════════════════
\section*{ASSIGNMENT 1 --- Full Solutions with Explanations}

\subsubsection*{Question 1}
\textit{The school of psychology known as functionalism, which emphasized the purposes of the mind's operations, was founded by:}
\medskip

\begin{tabular}{L{0.3cm} L{14.9cm}}
& A.\ Wilhelm Wundt \\
& B.\ John Watson \\
& \textbf{C.\ William James \checkmark} \\
& D.\ Edward Titchener \\
\end{tabular}

\vspace{0.3cm}
\noindent\textbf{Ans:} C — William James\\[0.2cm]
\textbf{Why this answer:} \\
Functionalism was founded by \textbf{William James}, who emphasized understanding mental processes in terms of their \textit{adaptive purpose} rather than their structure. Unlike structuralism, functionalism asked \textit{why} the mind operates the way it does, focusing on the functions and purposes of mental activity.
\vspace{0.4cm}

\subsubsection*{Question 2}
\textit{The fact that you know you have seen the face before but cannot remember the name of the person, illustrates the cognitive process of:}
\medskip

\begin{tabular}{L{0.3cm} L{14.9cm}}
& A.\ Perception \\
& B.\ Attention \\
& \textbf{C.\ Recognition \checkmark} \\
& D.\ Recall \\
\end{tabular}

\vspace{0.3cm}
\noindent\textbf{Ans:} C — Recognition\\[0.2cm]
\textbf{Why this answer:} \\
\textbf{Recognition} involves identifying previously encountered information without necessarily retrieving full contextual details. In this case, familiarity with the face is triggered, but the name — a contextually associated detail — cannot be retrieved. This is distinct from \textbf{recall}, which requires actively retrieving information from memory without a present cue.
\vspace{0.4cm}

\subsubsection*{Question 3}
\textit{The term \*\*\_\*\* refers to the relevance of the research to the "real world."}
\medskip

\begin{tabular}{L{0.3cm} L{14.9cm}}
& A.\ Internal validity \\
& B.\ Evolutionary validity \\
& \textbf{C.\ Ecological validity \checkmark} \\
& D.\ Neural validity \\
\end{tabular}

\vspace{0.3cm}
\noindent\textbf{Ans:} C — Ecological Validity\\[0.2cm]
\textbf{Why this answer:} \\
\textbf{Ecological validity} refers to the degree to which research findings can be generalized to real-world settings outside of the laboratory. High ecological validity means the study's conditions closely mirror natural, everyday situations. This concept is especially important in cognitive psychology, where lab tasks may not always reflect how cognition works in daily life.
\vspace{0.4cm}
\newpage
\subsubsection*{Question 4}
\textit{\_\_\_ challenged the assumption that mental events were beyond the realm of scientific study.}
\medskip

\begin{tabular}{L{0.3cm} L{14.9cm}}
& \textbf{A.\ Cognitive revolution \checkmark} \\
& B.\ Behaviorist rebellion \\
& C.\ Human factor movements \\
& D.\ Universal grammar \\
\end{tabular}

\vspace{0.3cm}
\noindent\textbf{Ans:} A — Cognitive Revolution\\[0.2cm]
\textbf{Why this answer:} \\
The \textbf{cognitive revolution} (circa 1950s–1960s) directly challenged behaviorism's dismissal of mental events. It re-established the legitimacy of studying internal mental processes — such as memory, attention, and problem-solving — using scientific methods. Key contributors included Noam Chomsky, George Miller, and Ulric Neisser.
\vspace{0.4cm}

\subsubsection*{Question 5}
\textit{Since the 1970s, various techniques of \*\*\_\*\* have allowed us to construct pictures of the anatomy and functioning of intact brains.}
\medskip

\begin{tabular}{L{0.3cm} L{14.9cm}}
& A.\ Genetic engineering \\
& \textbf{B.\ Brain imaging \checkmark} \\
& C.\ Neurosurgery \\
& D.\ Brain lesioning \\
\end{tabular}

\vspace{0.3cm}
\noindent\textbf{Ans:} B — Brain Imaging\\[0.2cm]
\textbf{Why this answer:} \\
\textbf{Brain imaging techniques} developed since the 1970s — including CT scans, PET scans, fMRI, and EEG — allow researchers to examine both the structure and functioning of the intact human brain \textit{non-invasively}. These tools have been transformative for cognitive neuroscience and cognitive psychology.
\vspace{0.4cm}

\subsubsection*{Question 6}
\textit{Historians date the founding of scientific psychology to the 1879 laboratory of:}
\medskip

\begin{tabular}{L{0.3cm} L{14.9cm}}
& A.\ William James \\
& \textbf{B.\ Wilhelm Wundt \checkmark} \\
& C.\ John Locke \\
& D.\ Edward Titchener \\
\end{tabular}

\vspace{0.3cm}
\noindent\textbf{Ans:} B — Wilhelm Wundt\\[0.2cm]
\textbf{Why this answer:} \\
\textbf{Wilhelm Wundt} established the first formal experimental psychology laboratory in Leipzig, Germany in \textbf{1879}, marking the birth of psychology as a distinct scientific discipline. Wundt's approach, known as \textbf{structuralism}, used introspection to analyze the basic elements of conscious experience.
\vspace{0.4cm}
\newpage

\subsubsection*{Question 7}
\textit{The term "limited capacity processors" suggests that:}
\medskip

\begin{tabular}{L{0.3cm} L{14.9cm}}
& A.\ Humans have limited memory storage \\
& \textbf{B.\ Human beings can only do so many things at once \checkmark} \\
& C.\ Neurons can only fire at a certain rate \\
& D.\ Processing in the brain occurs in parallel \\
\end{tabular}

\vspace{0.3cm}
\noindent\textbf{Ans:} B — Human beings can only do so many things at once\\[0.2cm]
\textbf{Why this answer:} \\
The concept of \textbf{limited-capacity processors} in cognitive psychology means that human cognitive resources — including attention, working memory, and processing — are finite. This limits how many tasks or how much information can be handled simultaneously. This principle is central to the \textbf{information-processing approach} to cognition.
\vspace{0.4cm}

\subsubsection*{Question 8}
\textit{This psychologist sought to describe the intellectual structures underlying cognitive experience at different developmental points through an approach he called genetic epistemology:}
\medskip

\begin{tabular}{L{0.3cm} L{14.9cm}}
& A.\ Galton \\
& B.\ Skinner \\
& C.\ Wundt \\
& \textbf{D.\ Piaget \checkmark} \\
\end{tabular}

\vspace{0.3cm}
\noindent\textbf{Ans:} D — Piaget\\[0.2cm]
\textbf{Why this answer:} \\
\textbf{Jean Piaget} developed \textbf{genetic epistemology} — a framework to explain how knowledge and cognitive structures develop across different stages of childhood (sensorimotor, preoperational, concrete operational, and formal operational). His work was foundational to developmental and cognitive psychology.
\vspace{0.4cm}

\subsubsection*{Question 9}
\textit{Clinical interviews are to introspection as \_\_ are to \_\_\_\_.}
\medskip

\begin{tabular}{L{0.3cm} L{14.9cm}}
& A.\ Naturalistic observations; experiments \\
& B.\ Naturalistic observations; quasi experiments \\
& \textbf{C.\ Controlled observations; naturalistic observations \checkmark} \\
& D.\ Controlled observations; experiments \\
\end{tabular}

\vspace{0.3cm}
\noindent\textbf{Ans:} C — Controlled observations; naturalistic observations\\[0.2cm]
\textbf{Why this answer:} \\
This is an \textbf{analogy question}. Clinical interviews are a \textit{structured/controlled} version of introspection, just as \textbf{controlled observations} are a more structured version of \textbf{naturalistic observations}. The analogy highlights the relationship between less controlled, more naturalistic methods and their structured counterparts.
\vspace{0.4cm}

\subsubsection*{Question 10}
\textit{Which of the following is a basic assumption of the connectionist approach?}
\medskip

\begin{tabular}{L{0.3cm} L{14.9cm}}
& A.\ Serial processing \\
& B.\ Multiple stores where information is kept throughout processing \\
& \textbf{C.\ Networks of connections among simple processing units \checkmark} \\
& D.\ A central processor that directs the flow of information \\
\end{tabular}

\vspace{0.3cm}
\noindent\textbf{Ans:} C — Networks of connections among simple processing units\\[0.2cm]
\textbf{Why this answer:} \\
The \textbf{connectionist (or parallel distributed processing) approach} assumes that cognition emerges from vast \textbf{networks of interconnected simple units} that process information simultaneously. This differs from the classical information-processing model, which assumes a central processor and serial processing. Connectionism draws inspiration from the neural architecture of the brain.
\vspace{0.4cm}

\medskip
\subsection*{Assignment Quick Answer Summary}

\begin{center}
\renewcommand{\arraystretch}{1.4}
\begin{tabular}{|L{0.8cm}|L{6.2cm}|L{8.2cm}|}
\hline
\bb{Q\#} & \bb{Topic} & \bb{Correct Answer} \\
\hline
1 & Founding of Functionalism & C — William James \\
2 & Recognition vs. Recall & C — Recognition \\
3 & Research Validity & C — Ecological Validity \\
4 & Challenging Behaviorism & A — Cognitive Revolution \\
5 & Brain Study Techniques & B — Brain Imaging \\
6 & Birth of Scientific Psychology & B — Wilhelm Wundt \\
7 & Limited Capacity Processors & B — Can only do so many things at once \\
8 & Genetic Epistemology & D — Piaget \\
9 & Research Methods Analogy & C — Controlled; Naturalistic Observations \\
10 & Connectionist Approach & C — Networks of simple processing units \\
\hline
\end{tabular}
\end{center}

\medskip
\end{document}
